% Options for packages loaded elsewhere
\PassOptionsToPackage{unicode}{hyperref}
\PassOptionsToPackage{hyphens}{url}
\documentclass[
]{article}
\usepackage{xcolor}
\usepackage{amsmath,amssymb}
\setcounter{secnumdepth}{-\maxdimen} % remove section numbering
\usepackage{iftex}
\ifPDFTeX \usepackage[T1]{fontenc} \usepackage[utf8]{inputenc} \usepackage{textcomp} % provide euro and other symbols
\else % if luatex or xetex \usepackage{unicode-math} % this also loads fontspec \defaultfontfeatures{Scale=MatchLowercase} \defaultfontfeatures[\rmfamily]{Ligatures=TeX,Scale=1}
\fi
\usepackage{lmodern}
\ifPDFTeX\else % xetex/luatex font selection
\fi
% Use upquote if available, for straight quotes in verbatim environments
\IfFileExists{upquote.sty}{\usepackage{upquote}}{}
\IfFileExists{microtype.sty}{% use microtype if available \usepackage[]{microtype} \UseMicrotypeSet[protrusion]{basicmath} % disable protrusion for tt fonts
}{}
\makeatletter
\@ifundefined{KOMAClassName}{% if non-KOMA class \IfFileExists{parskip.sty}{% \usepackage{parskip} }{% else \setlength{\parindent}{0pt} \setlength{\parskip}{6pt plus 2pt minus 1pt}}
}{% if KOMA class \KOMAoptions{parskip=half}}
\makeatother
\usepackage{graphicx}
\makeatletter
\newsavebox\pandoc@box
\newcommand*\pandocbounded[1]{% scales image to fit in text height/width \sbox\pandoc@box{#1}% \Gscale@div\@tempa{\textheight}{\dimexpr\ht\pandoc@box+\dp\pandoc@box\relax}% \Gscale@div\@tempb{\linewidth}{\wd\pandoc@box}% \ifdim\@tempb\p@<\@tempa\p@\let\@tempa\@tempb\fi% select the smaller of both \ifdim\@tempa\p@<\p@\scalebox{\@tempa}{\usebox\pandoc@box}% \else\usebox{\pandoc@box}% \fi%
}
% Set default figure placement to htbp
\def\fps@figure{htbp}
\makeatother
\setlength{\emergencystretch}{3em} % prevent overfull lines
\providecommand{\tightlist}{% \setlength{\itemsep}{0pt}\setlength{\parskip}{0pt}}
\usepackage{bookmark}
\IfFileExists{xurl.sty}{\usepackage{xurl}}{} % add URL line breaks if available
\urlstyle{same}
\hypersetup{ hidelinks, pdfcreator={LaTeX via pandoc}} \author{}
\date{} \begin{document} % KAPAK SAYFASI
\begin{titlepage}
\centering \vspace*{0.5cm} \includegraphics[width=3in]{media/media/image5.png} \vspace{0.5cm} {\LARGE\bfseries CNG 465\par} \vspace{0.5cm} {\Large INTRODUCTION TO BIOINFORMATICS\par} \vspace{1cm} {\large Fall 25-26\par} \vspace{0.3cm} {\large Semester Project Report\par} \vspace{0.3cm} {\large \par} \vfill {\large \par
\vspace{0.2cm} \par
\vspace{0.2cm} \par
\vspace{0.2cm}
Zanyar Erkozan 2584985\par
} \vspace{1cm} \end{titlepage} \newpage % ABSTRACT
\section{ABSTRACT}\label{abstract} The primary goal of the project was to design and develop a family tree
and to design algorithms implementing family relationships based on this
design. It is a model of individuals, relationships of the parent and
child, and relationships of marriage. The project was written in Python programming language using object
orientation. The basic classes were deployed to handle persons,
gender-specific persons, partnerships and traversal nodes. Then
different functions were built on the developed structure based on the
relationships within the families. This includes parents, kids,
ancestors, descendants and other relationships through various graph
methods by traversals. The developed system makes the traversal of the family relationship in a
systematic method with the help of breadth-first search algorithm. To
ensure the correct values of the developed system has been confirmed
through various simulations performed on different family structures.
The simulation results indicate that the developed system is capable of
displaying the family relationships in a proper manner and performing
genealogical queries effectively. \newpage % TABLE OF CONTENTS
\section{Table of Contents}\label{table-of-contents} \hyperref[abstract]{\textbf{1. ABSTRACT 2}} \hyperref[table-of-contents]{\textbf{Table of Contents 3}} \hyperref[table-of-figures]{\textbf{Table of Figures 4}} \hyperref[introduction]{\textbf{2. INTRODUCTION 5}} \hyperref[contribution]{\textbf{3.} \textbf{CONTRIBUTION} \textbf{5}} \hyperref[problem-statement]{\textbf{4. PROBLEM STATEMENT 6}} \hyperref[solution]{\textbf{5.} \textbf{SOLUTION} \textbf{7}} \begin{quote}
\hyperref[partnership-class]{1) Partnership Class 7} \hyperref[str__-function]{A) \_\_str\_\_ Function 8} \hyperref[repr__-function]{B) \_\_repr\_\_ Function 9} \hyperref[is_parent-function]{C) is\_parent Function 10} \hyperref[is_mother-function]{D) is\_mother Function 11} \hyperref[is_father-function]{E) is\_father Function 12} \hyperref[add_parent-function]{F) add\_parent Function 13} \hyperref[add_child-function]{G) add\_child Function 14} \hyperref[add_children-function]{H) add\_children Function 15} \hyperref[traverse-function]{I) traverse Function 15} \hyperref[person-class]{2) Person Class 21} \hyperref[str__-function-1]{A) \_\_str\_\_ Function 22} \hyperref[eq__-function]{B) \_\_eq\_\_ Function 23} \hyperref[ancestor-function]{C) ancestor Function 24} \hyperref[descendant-function]{D) descendant Function 26}
\end{quote} \hyperref[results]{\textbf{6. RESULTS 28}} \hyperref[data-flow-diagram]{\textbf{7.DATA FLOW DIAGRAM 29}} \hyperref[conclusion]{\textbf{8.} \textbf{CONCLUSION} \textbf{30}} \hyperref[references]{\textbf{9. REFERENCES 30}} \newpage % TABLE OF FIGURES
\section{Table of Figures}\label{table-of-figures} \begin{quote}
\hyperref[figure-1-__str__-function-at-partnership-class]{Figure 1
\_\_str\_\_ Function at Partnership class 8} \hyperref[figure-2-__repr__-function-at-partnership-class]{Figure 2
\_\_repr\_\_ Function at Partnership class 9} \hyperref[figure-3-is_parent-function-at-partnership-class]{Figure 3
is\_parent Function at Partnership class 10} \hyperref[figure-4-is_mother-function-at-partnerships-class]{Figure 4
is\_mother Function at Partnerships class 11} \hyperref[figure-5-is_father-function-at-partnership-class]{Figure 5
is\_father Function at Partnership class 12} \hyperref[figure-6-add_parent-function-at-partnership-class]{Figure 6
add\_parent Function at Partnership class 13} \hyperref[figure-7-add_child-function-at-partnership-class]{Figure 7
add\_child Function at Partnership class 14} \hyperref[figure-8-add_children-function-at-partnership-class]{Figure 8
add\_children Function at Partnership class 15} \hyperref[figure-9-traverse-function-at-partnership-class]{Figure 9
traverse Function at Partnership class 18} \hyperref[figure-10-__str__-function-at-person-class]{Figure 10
\_\_str\_\_ Function at Person class 22} \hyperref[figure-11-__eq__-function-at-person-class]{Figure 11
\_\_eq\_\_ Function at Person class 23} \hyperref[figure-12-ancestor-function-at-person-class]{Figure 12
ancestor Function at Person class 24} \hyperref[figure-13-descendants-function-at-person-class]{Figure 13
descendants Function at Person class 26} \hyperref[figure-14-test-results]{Figure 14 Test Results 28} \hyperref[figure-15-data-flow-diagram]{Figure 15 Data Flow Diagram 29}
\end{quote} \newpage % INTRODUCTION
\section{INTRODUCTION}\label{introduction} The objective of this Project was to develop a family tree system with
the capability of organizing family relations in an organized way. The
use of family trees are common in applications concerning bioinformatics
as there is a need to display the relationships of the various people
efficiently. In the Project a family relationship network, a graph structure was
created using object oriented programming with Python. Members of a
family and their relationship\\
information including gender relationships and parent and child
relationships were\\
clearly represented using objects. Functions were also created using a
graph structure including parents, children and other relationships of a
person depending on the graph traversal methods. Not only was the aim of the project to store family data but also to
analyze the relations systematically which means traversal algorithms.
By doing so the system is able to analyze both the design process
implementations and testing the code of the developed system. \newpage % CONTRIBUTION
\section{CONTRIBUTION}\label{contribution} In this Project, () and ()
was involved mainly in coding implementation and testing. Drawing the Project's logic algorithm flowchart Zanyar Erkozan(2584985)
was involved. In the Report Writing, translating functions to LaTeX format, describing
the system design and functionality and documenting the system, () and () was involved. \newpage % PROBLEM STATEMENT
\section{PROBLEM STATEMENT}\label{problem-statement} The primary problem that is handled in the project is the fact that
there is no efficient software to handle relationships between families
and to traverse through the relationships represented in the software.
This is because the easiest way to represent these relationships between
the families will only involve the parents or the children, but it is
not effective when relationships are to be queried both up and downward. Another problem is related to ensuring consistency while traversing the
family structure. When traversing the tree, it is necessary that the
persons being traversed are not traversed more than once. To add to
this, it is required for it to process the relationships in the correct
way in the same framework for the relationships like ancestors,
descendants, siblings and cousins. Thus, the problem statement can be conceptualized as encapsulating the
design of a family tree system such that: correctly reflects persons,
sexes, partnerships and relationships of parents and children, provides
support for systematic traversal of the family structure, simplifies
looking into multiple family relationships in one common manner and is
also scalable for additional analysis tasks. This project proposes the resolution of these problems by modeling the
family relations using a representation involving graphs and by building
functions using traversals to examine said relationships. \newpage % SOLUTION
\section{SOLUTION}\label{solution} To solve the problem that was defined a graph model of the family
structure was developed using object oriented programming techniques and
Python programming.\\
The rationale was that people and their relationships could be created
explicitly and these relationships could then be evaluated using formal
graph traversal techniques. In this solution persons, partnerships, parent and child relationships
were modeled using distinct classes. Thus this was beneficial in making
sure that the system is able to differentiate between persons and their
relationships. According to this methods were developed that allow
traversal through the family tree. Every key function in the system traces an algorithm sequence clearly.
Thus in this respect the algorithms used in the function implementation
will be represented in the formal algorithm description style by
applying the use of the LaTeX format as used in the standard algorithm
definition in formal literature. In the following subsections, the primary algorithms employed in the
system are explained. In each instance, the LaTeX expression
representing the algorithm is provided first, and then an explanation is
given of its application and functions. This is done to help the reader
comprehend both the high-level view and the working of the developed
solution. \subsection{Partnership Class}\label{partnership-class} A partnership is a representation of the parent pair and their offspring
in the family tree. It can also be perceived as a marriage/partnership
node which joins the pair of parents to their offspring in one unit. This gives room to model family relationships in a straightforward
manner .Within this class, the parents attribute holds the mother and
father of a relationship, and the children attribute holds a list of all
the children within a relationship. Along with these relationships, this
class is linked to the person class by back-references. The child holds
a reference to a relationship (child\_backref), and a parent holds a
reference to all relationships within which it is a parent
(partnership\_backrefs). This creates a relationship tree that can be
traversed forward and backward efficiently. Further into the phase of the project, various methods in the
partnership class were enabled through the instructor in the context of
the project framework. These methods are \_\_str\_\_, \_\_repr\_\_,
is\_parent, is\_mother, is\_father, add\_parent, add\_child,
add\_children, and the `traverse` function. These methods were each
incorporated to perform distinct duties in the context of family
relationships, to be clarified in the coming sections. \subsubsection{\_\_str\_\_ Function}\label{str__-function} \_\_str\_\_ function defines how a partnership object is represented as
a string. When the object is printed or converted to a string, it
returns the partnership's name if it exists. If there is no name
whatsoever it returns the default value "Partnership". \includegraphics[width=6.26772in,height=1.86111in,alt={Figure 1 Partnership class \_\_str\_\_ function }]{media/media/image12.png} \paragraph{\texorpdfstring{ Figure 1 \_\_str\_\_ Function at Partnership
class}{ Figure 1 \_\_str\_\_ Function at Partnership class}}\label{figure-1-__str__-function-at-partnership-class} Line 1: Defines the \_\_str\_\_ method for a partnership object. This
method decides what string will be shown when the object is printed. Line 2: Checks whether the partnership has a valid name. If name does
exists and is not empty, it will be used. Line 3: Returns the partnership name as the string representation. Line 4: The alternative case when the partnership has no name. Line 5: Returns the default string "Partnership" so the object still has
a readable output. \subsubsection{\_\_repr\_\_ Function}\label{repr__-function} The \_\_repr\_\_ function is acting like the \_\_str\_\_ function in
this project, which means it defines how a partnership object is
represented whenever it would appear in the console or when embedded
into another data structure, such as a list. If the partnership has a
name, return the name; else use some default string. \includegraphics[width=6.26772in,height=1.98611in]{media/media/image16.png} \paragraph{\texorpdfstring{ \emph{Figure 2 \_\_repr\_\_ Function at
Partnership
class}}{ Figure 2 \_\_repr\_\_ Function at Partnership class}}\label{figure-2-__repr__-function-at-partnership-class} Line 1: Definition of the \_\_repr\_\_ method for a partnership object,
which controls its representation. Line 2: Checks whether the partnership has a valid name. Line 3: Returns the name of the partnership if it exists. Line 4: The alternative case when the partnership has no name. Line 5: Returns the default string "Partnership" to ensure a readable
output. \subsubsection{is\_parent Function}\label{is_parent-function} The is\_parrent function checks if the individual is parent or not. The
function returns the value True if the individual is mother or father.
If the individual is not a parent the code automatically returns the
value False for code stability in the future usage. \includegraphics[width=6.26772in,height=2.38889in]{media/media/image18.png} \paragraph{\texorpdfstring{ \emph{Figure 3 is\_parent Function at
Partnership
class}}{ Figure 3 is\_parent Function at Partnership class}}\label{figure-3-is_parent-function-at-partnership-class} Line 1: Definition of the is\_parrent method for a partnership object,
which checks the individual is parent or not. Line 2: Checks whether the partnership's mother is the individual's
mother . Line 3: Returns true if it is correct. Line 4: Checks whether the partnership's father is the individual's
father. Line 5: Returns true if it is correct. Line 6: Returns the default answer as False to ensure stability. \subsubsection{is\_mother Function}\label{is_mother-function} The is\_mother function checks if the individual is a mother or not. The
function returns the value True if the individual is a mother. If the
individual is not a mother the code automatically returns the value
False for code stability in the future usage. \includegraphics[width=6.26772in,height=1.75in]{media/media/image15.png} \paragraph{\texorpdfstring{ Figure 4 is\_mother Function at Partnerships
class}{ Figure 4 is\_mother Function at Partnerships class}}\label{figure-4-is_mother-function-at-partnerships-class} Line 1: Definition of the is\_mother method for a partnership object,
which checks the individual is mother or not. Line 2: Checks whether the partnership's mother is the individual's
mother . Line 3: Returns true if the individual is the mother. Line 4: Returns false if the individual is not a mother. \subsubsection{is\_father Function}\label{is_father-function} The is\_father function checks if the individual is father or not. The
function returns the value True if the individual is father. If the
individual is not a father, the code automatically returns the value
False for code stability in the future usage. \includegraphics[width=6.26772in,height=1.70833in]{media/media/image11.png} \paragraph{\texorpdfstring{ \emph{Figure 5 is\_father Function at
Partnership
class}}{ Figure 5 is\_father Function at Partnership class}}\label{figure-5-is_father-function-at-partnership-class} Line 1: Definition of the is\_father method for a partnership and
individual object, which checks the individual is father or not. Line 2: Checks whether the partnership's mother is the individual's
father. Line 3: Returns true if the individual is the father. Line 4: Returns true if the individual is not the father. \subsubsection{add\_parent Function}\label{add_parent-function} The add\_parrent function adds the partnership to the parent list. The
function checks the all of the parents and calls the add\_parent
function with the partnership and parent parameters. \includegraphics[width=6.26772in,height=2.55556in]{media/media/image9.png} \paragraph{\texorpdfstring{ \emph{Figure 6 add\_parent Function at
Partnership
class}}{ Figure 6 add\_parent Function at Partnership class}}\label{figure-6-add_parent-function-at-partnership-class} Line 1: Definition of the is\_parrent method for a partnership object,
which checks the individual is parent or not. Line 2: Checks whether the partnership's mother is the individual's
mother . Line 3: Returns true if it is correct. Line 4: Checks whether the partnership's father is the individual's
father. Line 5: Returns true if it is correct. Line 6: Returns the default answer as False to ensure stability. \subsubsection{add\_child Function}\label{add_child-function} The add\_child function's main task is to add a child to the
partnership. This is done by checking if the child does not exist in the
partnerships child. Then adds the child to the partnership. Then sets
the child's back reference to the partnership which is very important
because in the code we can go through child to partnership as well as
with the back reference we can go through partnership to child. \includegraphics[width=6.26772in,height=1.75in]{media/media/image17.png} \paragraph{\texorpdfstring{ Figure 7 add\_child Function at Partnership
class}{ Figure 7 add\_child Function at Partnership class}}\label{figure-7-add_child-function-at-partnership-class} Line 1: Defines the add\_child method with partnership and child as an
input object. Line 2: Checks whether the child is the partnership child or not. Line 3: If the child is not a child of the partnership, the code adds
the child to the partnership. Line 4: If the child is not a child of the partnership, we set the
child's back reference to the partnership. \subsubsection{add\_children Function}\label{add_children-function} The add\_childeren function's main task is to add multiple children to
the partnership. This is done by adding each child in the children list
by calling the add\_child function. \includegraphics[width=6.26772in,height=1.41667in]{media/media/image14.png} \paragraph{\texorpdfstring{ \emph{Figure 8 add\_children Function at
Partnership
class}}{ Figure 8 add\_children Function at Partnership class}}\label{figure-8-add_children-function-at-partnership-class} Line 1: Defines the add\_children method with partnership and child as
an input object. Line 2:For loop gets one child in the children one at a time. Line 3: Calls the add\_child function with the input parameters as
partnership and child. \subsubsection{traverse Function}\label{traverse-function} The traverse function is using breadth-first search traversal in the
family tree from a given person. The search passes through parents,
children, siblings, and relationships such as cousins. Each visited
person is recorded along with a level indicating their position relative
to each individual in the tree. The search also keeps record of the
individual with the lowest level encountered. \includegraphics[width=6.26772in,height=7.97222in]{media/media/image8.png} \includegraphics[width=6.26772in,height=7.88889in]{media/media/image1.png}\includegraphics[width=6.62188in,height=1.41022in]{media/media/image6.png} \paragraph{}\label{section} \paragraph{\texorpdfstring{ \emph{Figure 9 traverse Function at
Partnership
class}}{ Figure 9 traverse Function at Partnership class}}\label{figure-9-traverse-function-at-partnership-class} Line 1: Defines the traverse method with start\_person and start\_level
as an input. Line 2: Makes visited\_nodes to be an empty list. Line 3: Makes visited\_persons as an empty set. Line 4: Makes a queue with the start\_person and start\_level. Line 5: Adds start\_person to the visited\_persons. Line 6: Sets min\_level to +infinity. Line 7: Sets min\_level\_node to NULL. Line 8: While loop that checks if the queue is not empty. Line 9: Removes the first element from the queue. Line 10: Creates a visit\_node class vn with current\_person and level. Line 11: Adds vn to the visited\_nodes. Line 12: Checks that if the level is less than min\_level. Line 13: If level is less than min\_level, sets level to min\_level. Line 14: If level is less than min\_level, sets vn to the
min\_level\_node. Line 15: Checks if visit\_callback exists. Line 16: If visit\_callback exists, calls visit\_callback function with
current\_person and level. Line 17: Checks current person has any child\_backref. Line 18: If current person has child\_backref, sets current\_person's
child\_backref to partnership. Line 19: If current person has child\_backref, checks if mother exists
and not visited. Line 20: If current person has child\_backref and if the mother exists
and not visited, adds mother to visited person. Line 21: If current person has child\_backref and if the mother exists
and not visited, adds mother to the queue with the level of level-1. Line 22: If current person has child\_backref, checks if the father
exists and not visited. Line 23:If current person has child\_backref and checks if the father
exists and not visited, adds father to the visited\_persons. Line 24:If current person has child\_backref and checks if the father
exists and not visited, adds father to the queue with the level of
level-1. Line 25: Creates a loop for each partnership in current\_person's
partnership\_backrefs. Line 26: Creates a loop for each child in partnership.children while
maintaining the loop for each partnership in current\_person's
partnership\_backrefs. Line 27: Checks if a child is not visited as well as within a loop for
each child in partnership.children while maintaining the loop for each
partnership in current\_person's partnership\_backrefs. Line 28: If a child is not visited as well as within a loop for each
child in partnership.children while maintaining the loop for each
partnership in current\_person's partnership\_backrefs, adds child to
the visited\_persons Line 29: If a child is not visited as well as within a loop for each
child in partnership.children while maintaining the loop for each
partnership in current\_person's partnership\_backrefs, adds child to
the queue with the level of level -1. Line 30: Checks if a current\_person has a child\_backref Line 31: If a current\_person has a child\_backref, sets
current\_person's child\_backref to partnership. Line 32: If a current\_person has a child\_backref while maintaining the
loop for each partnership in current\_person's partnership\_backrefs,
makes a loop for each sibling in partnership's children. Line 33: If a current\_person has a child\_backref with a loop for each
sibling in partnership's children, checks if the sibling is not equal to
the current\_person and not visited. Line 34: If a current\_person has a child\_backref with a loop for each
sibling in partnership's children as well as with the sibling is not
equal to the current\_person and not visited, adds sibling to
visited\_person. Line 35: If a current\_person has a child\_backref with a loop for each
sibling in partnership's children as well as with the sibling is not
equal to the current\_person and not visited, adds sibling to the queue
with the level of level. Line 36: Checks if a current\_person has a child\_backref. Line 37: If a current\_person has a child\_backref, set
current\_person's child\_backref to partnership. Line 38: If a current\_person has a child\_backref, makes a loop for
each parent in mother and father. Line 39: If a current\_person has a child\_backref with a loop for each
parent in mother and father, checks if parent exists and has a
child\_backref. Line 40: If a current\_person has a child\_backref with a loop for each
parent in mother and father and parent exists and has a child\_backref,
sets parent's child\_backref to parent\_partnership. Line 41: If a current\_person has a child\_backref with a loop for each
parent in mother and father as well as if parent exists and has a
child\_backref, creates a loop for each parent\_sibling in
parent\_partnership's children. Line 42: If a current\_person has a child\_backref with a loop for each
parent in mother and father as well as if parent exists and has a
child\_backref while satisfying a loop for each parent\_sibling in
parent\_partnership's children, checks if parent\_sibling is not equal
to parent and not visited. Line 43: If a current\_person has a child\_backref with a loop for each
parent in mother and father as well as if parent exists and has a
child\_backref while satisfying a loop for each parent\_sibling in
parent\_partnership's children and if parent\_sibling is not equal to
parent and not visited, adds parent\_sibling to visited\_persons. Line 44: If a current\_person has a child\_backref with a loop for each
parent in mother and father as well as if parent exists and has a
child\_backref while satisfying a loop for each parent\_sibling in
parent\_partnership's children and if parent\_sibling is not equal to
parent and not visited, adds parent\_sibling to a queue with the level
of level -1. Line 45: If a current\_person has a child\_backref with a loop for each
parent in mother and father as well as if parent exists and has a
child\_backref while satisfying a loop for each parent\_sibling in
parent\_partnership's children, makes a loop for each partnership in
parent\_sibling. Line 46: If a current\_person has a child\_backref with a loop for each
parent in mother and father as well as if parent exists and has a
child\_backref while satisfying a loop for each parent\_sibling in
parent\_partnership's children, makes a loop for each partnership in
parent\_sibling. Line 47: If a current\_person has a child\_backref with a loop for each
parent in mother and father as well as if parent exists and has a
child\_backref while satisfying a loop for each parent\_sibling in
parent\_partnership's children and fulfills the loop for each
partnership in parent\_sibling, checks if cousin is not visited. Line 48: If a current\_person has a child\_backref with a loop for each
parent in mother and father as well as if parent exists and has a
child\_backref while satisfying a loop for each parent\_sibling in
parent\_partnership's children and fulfills the loop for each
partnership in parent\_sibling and if cousin is not visited, adds cousin
to visited\_persons. Line 49: If a current\_person has a child\_backref with a loop for each
parent in mother and father as well as if parent exists and has a
child\_backref while satisfying a loop for each parent\_sibling in
parent\_partnership's children and fulfills the loop for each
partnership in parent\_sibling and if cousin is not visited, adds cousin
to a queue with the level of level. Line 50: Checks if level sorting is enabled. Line 51: If level sorting is enabled, sorts visited\_nodes by level. Line 52: Returns visited\_nodes and min\_level\_node. \subsection{2) Person Class}\label{person-class} The Person class represents an individual in the family tree. Each
person can be assigned a name and an optional DNA code,and can solve
extra information using a flexible properties dictionary. The class is
designed to support family relationship queries by linking individuals
to partnerships through back-referances.\\
\strut \\
A person might be connected to the family tree in to different roles.
Firstly a person can be a child of a certain partnership using
child\_bacref which points to the partnership where the person appears
as child. Secondly one can be parent in one or more partnerships which
is realized by partnership\_backrefs storing a list of partnerships
where this person is registered as parent. These links make it possible
to move upwards to parents and downwards to children when navigating a
family tree. In the project framework, the Person class includes a number utility and
relationship related methods. While some have been implemented by our
team with the required functionality, others were provided as part of
the base structure of this project. In what follows we focus on these
methods provided by our group. The methods implemented by our team and
explained in this reports are: \_\_str\_\_, \_\_eq\_\_, ancestor and
descendant. \subsubsection{\_\_str\_\_ Function}\label{str__-function-1} The `\_\_str\_\_' function returns a string representation of a person
id a person has a name thats what returned otherwise "Unknown" is set as
a default output. \includegraphics[width=6.26772in,height=2.08333in]{media/media/image7.png} \paragraph{\texorpdfstring{ Figure 10 \_\_str\_\_ Function at Person
class}{ Figure 10 \_\_str\_\_ Function at Person class}}\label{figure-10-__str__-function-at-person-class} Line 1: Defines the method \_\_str\_\_ which decides how the person
object will be represented as a string or printed out. Line 2: It checks if the individual has a valid name. Line 3: This returns the person\textquotesingle s name if it is
available. Line 4: Starts the alternative case if the name does not exits. Line 5:This is a return statement with the default value "UNknown" to
make sure the object is always printable. \subsubsection{\_\_eq\_\_ Function}\label{eq__-function} The ``\_\_eq\_\_'' method is used to compare the equality of the objects
of the persons class. Two persons are said to be equal if both are of
the same class and have the same name. The method will return False if
the object passed is not a person object. \includegraphics[width=6.26772in,height=2.80556in]{media/media/image10.png} \paragraph{\texorpdfstring{ Figure 11 \_\_eq\_\_ Function at Person
class}{ Figure 11 \_\_eq\_\_ Function at Person class}}\label{figure-11-__eq__-function-at-person-class} Line 1: Defines the ``\_\_eq\_\_'' method for comparing the current
person with another object. Line 2:Checks whether the other object is an instanceof the person
class. Line 3:If the other object is an instanceof the person class, this line
will return false if the passed object is not a person. Line 4: Checks if the names of both person objects are equal. Line 5: If the names of both person objects are equal, returns true. Line 6:Checks if the names of both person objects are not equal. Line 7:If the names of both person objects are not equal,It will return
false if the names don't match. \subsubsection{ancestor Function}\label{ancestor-function} The ancestor function returns a list that includes all ancestors of a
dedicated person traversing upward through the family tree. The parent
links are traversed based on child\_backref pointer as well as mother
and father at different levels. The visited set used to avert returning
a person or processing it to avoid loops. \includegraphics[width=6.26772in,height=6.25in]{media/media/image4.png} \paragraph{Figure 12 ancestor Function at Person
class}\label{figure-12-ancestor-function-at-person-class} Line 1: Begins the ancestor method for a person Line 2: Initializes ancestors\_list as an empty list. Line 3:Initializes visited as an empty set. Line 4:Defines recursive function as collect function with an input of
current\_person. Line 5:Defines recursive function as collect function with an input of
current\_person and checks if current\_person is in visited. Line 6:Defines recursive function as collect function with an input of
current\_person and adds current\_person to visited. Line 7:Defines recursive function as collect function with an input of
current\_person and if current\_person is in visited, returns current
person. Line 8:Defines recursive function as collect function with an input of
current\_person and checks if current\_person's child\_backref is not
null value. Line 9:Defines recursive function as collect function with an input of
current\_person and if current\_person's child\_backref is not null
value, sets current person's child\_backref to partnership. Line 10:Defines recursive function as collect function with an input of
current\_person and if current\_person's child\_backref is not null
value, checks if partnership's parent's mother is not null and mother is
not visited. Line 11: Defines recursive function as collect function with an input of
current\_person and if current\_person's child\_backref is not null
value and iif partnership's parent's mother is not null, checks mother
is not visited. Line 12:Defines recursive function as collect function with an input of
current\_person and if current\_person's child\_backref is not null
value and if partnership's parent's mother is not null and mother is not
visited, adds mother to anchestors\_list. Line 13:Defines recursive function as collect function with an input of
current\_person and if current\_person's child\_backref is not null
value and if partnership's parent's mother is not null and mother is not
visited, calls collect function with the input of mother. Line 14: Defines recursive function as collect function with an input of
current\_person and if current\_person's child\_backref is not null
value, checks if partnership's parents's father is not null and father
is not in visited. Line 15: Defines recursive function as collect function with an input of
current\_person and if current\_person's child\_backref is not null
value and if partnership's parents's father is not null, checks father
is not in visited. Line 16:Defines recursive function as collect function with an input of
current\_person and if current\_person's child\_backref is not null
value and if partnership's parents's father is not null and father is
not in visited, adds father to ancestors\_list. Line 17:Defines recursive function as collect function with an input of
current\_person and if current\_person's child\_backref is not null
value and if partnership's parents's father is not null and father is
not in visited, calls collect function with the input of father. Line 18:Calls collect function with the input of person. Line 19:Returns the ancestors\_list. \subsubsection{descendant Function}\label{descendant-function} The descendent method on the other hand traverses down a family tree to
return all descendants of a person. The method iterates over all pairing
that involve a person as a parent and append all their child to a set of
descendants. The method employs a visited set to ensure that a person is
not traversed more than once. The method is recursive meaning that it
calls it self to move down generations. \begin{quote}
\includegraphics[width=6.26772in,height=5in]{media/media/image2.png}
\end{quote} \paragraph{\texorpdfstring{ \emph{Figure 13 descendants Function at
Person
class}}{ Figure 13 descendants Function at Person class}}\label{figure-13-descendants-function-at-person-class} Line 1: Begins the descendants method for a person Line 2: Initializes descentants\_list as an empty list. Line 3:Initializes visited as an empty set. Line 4:Defines recursive function as collect function with an input of
current\_person. Line 5:Defines recursive function as collect function with an input of
current\_person and checks if current\_person is in visited. Line 6:Defines recursive function as collect function with an input of
current\_person and adds current\_person to visited. Line 7:Defines recursive function as collect function with an input of
current\_person and if current\_person is in visited, returns current
person. Line 8:Defines recursive function as collect function with an input of
current\_person and makes a loop for each partnership in
current\_person's partnership\_backrefs. Line 9:Defines recursive function as collect function with an input of
current\_person while satisfying a loop for each partnership in
current\_person's partnership\_backrefs, creates a loop for each child
in partnership's children. Line 10:Defines recursive function as collect function with an input of
current\_person while satisfying a loop for each partnership in
current\_person's partnership\_backrefs and fulfills a loop for each
child in partnership's children, checks if the child is not in visited. Line 11:Defines recursive function as collect function with an input of
current\_person while satisfying a loop for each partnership in
current\_person's partnership\_backrefs and fulfills a loop for each
child in partnership's children and if the child is not in visited, adds
child to descendants\_list. Line 12:Defines recursive function as collect function with an input of
current\_person while satisfying a loop for each partnership in
current\_person's partnership\_backrefs and fulfills a loop for each
child in partnership's children and if the child is not in visited,
calls collect function with the input of child. Line 13: Calls collect function with the input of person. Line 14: Returns descendats\_list. \newpage % RESULTS
\section{6. RESULTS}\label{results} The system was tested with a test: The test case is called as ``kinship
relation tests'' which has 7 tests within it:
``test\_unrelated\_xx16\_xy24'', which is a test for two people if they
are related or not. ``test\_unrelated\_xx19\_xx18'', which is a test for
two people if they are related or not.
``test\_first\_cousins\_xy22\_xx18'', which is a test for two people if
they are first cousins or not.
``test\_first\_cousins\_once\_removed\_xy24\_xx18'', which is a test for
two people if they are first cousins who are removed or not.
``test\_greatuncle\_greatnephew\_xy24\_xy15'', which is a test for two
people if they are greatnephew/greatuncle or not.
``test\_uncle\_niece\_xx46\_xy15'' which is a test for two people if
they are uncle/niece or not. and ``test\_no\_such\_person\_error'' which
is a test for a person if person is real or not. In Figure 14, after
completing the ``kinship tests'', it can be observed that all tests are
completed without any errors which means that the first test's result is
that two people are unrelated. For the second test's result is that two
people are unrelated. For the third test's results, two people are first
cousins. For the fourth test's result, two people are first cousins are
once removed. For the fifth test's result two people are greatuncle or
greatnephew. For the sixth test's results, two people are uncle or
niece. For the seventh test's result, the person is a real person. After this step the program runs a simulation and gives the results
called ``Relationship Analysis Results''. In these result it can be
observed that person T19\_F1\_XX16 and T19\_F1\_XY24 are unrelated,
person T19\_F1\_XX19 and T19\_F1\_XX18 are unrelated, person
T19\_F1\_XY22 and T19\_F1\_XX18 are first\_cousins, person T19\_F1\_XY24
and T19\_F1\_XX18 are first\_cousins\_once\_removed, person
T19\_F1\_XY24 and T19\_F1\_XY15 are greatuncle\_greatnephew, person
T19\_F1\_XX46 and T19\_F1\_XY15 are uncle\_niece. After these results,
we observe that the program runs successfully. \includegraphics[width=6.30729in,height=3.44525in]{media/media/image13.png} \paragraph{Figure 14 Test Results}\label{figure-14-test-results} \newpage % DATA FLOW DIAGRAM
\section{7.DATA FLOW DIAGRAM}\label{data-flow-diagram} \includegraphics[width=6.26772in,height=6.75in]{media/media/image3.png} \paragraph{\texorpdfstring{ Figure 15 Data Flow
Diagram}{ Figure 15 Data Flow Diagram}}\label{figure-15-data-flow-diagram} Figure 15 represents the digitally drawn Data Flow Diagram of the
system. \newpage % CONCLUSION
\section{8. CONCLUSION}\label{conclusion} In conclusion, The objective of this project was to develop a basic
family tree structure with analysis of knipship relations using Python.
With using both person and partnership classes a basic structure was
created to make it possible two way traversal in parent child
relationship with using back references. Multiple functions were
implemented to make it possible this structure including functions for
identifying ancestors and descendants and tracing paths to visited
people.\\
\strut \\
Upon this framework the logic for performing a knipship relationship was
built using lowest common ancestor search and distance calculations. The
system was able to detect basic relationships between cousins, uncles
etc. From both the example script and the unit tests it was confirmed
that everything is working flawlessly based on the given family
arrangement. All test are passing which validates that the system is
also able to calculate different relationships as well as handle with
errors by throwing exceptions\\
\strut \\
Although the project deals with the most general cases as outlined by
the project tasks there is room for including the system to involve even
more complex family relationships. Still the completion of the task
meets the requirements of the project and serves as the initial stage of
analyzing kinship by means of the symbolic family tree. \newpage % REFERENCES
\section{9. REFERENCES}\label{references} N/A \end{document}